\numberlesssection{ЗАКЛЮЧЕНИЕ}

В результате выполнения курсовой работы была достигнута поставленная цель: разработано программное обеспечение для интерактивного конструирования и визуализации кинематики шестереночных механизмов в трехмерном пространстве.

В ходе работы были решены следующие задачи:
\begin{enumerate}
	\item Выполнен анализ предметной области, в ходе которого были рассмотрены математические основы построения трехмерных изображений и методы моделирования механических связей.
	\item Формализовано описание объектов сцены: определена геометрия зубчатых колес, осей и основания. Для позиционирования элементов механизма выбрана и математически описана осевая система координат гексагональной сетки, что позволило обеспечить корректное размещение сцепленных шестерен.
	\item Обоснован выбор и осуществлена программная реализация базовых алгоритмов машинной графики. Для удаления невидимых поверхностей применен алгоритм Z-буфера, обеспечивающий корректное отображение пересекающейся геометрии. Реалистичность изображения достигнута за счет использования модели освещения Фонга и алгоритма планарных теней.
	\item Разработан алгоритм расчета кинематики, основанный на представлении механизма в виде графа и его обходе в ширину. Алгоритм успешно вычисляет передаточные отношения, угловые скорости и направления вращения для всей цепи, а также детектирует ситуации механического заклинивания.
	\item Спроектирована модульная структура программного обеспечения, выделяющая слои логики (сцена, кинематика) и представления (отрисовка, интерфейс). Приложение реализовано на языке C++ (стандарт C++20) с использованием библиотеки Qt.
	\item Проведено исследование разработанного программного комплекса. Результаты показали, что время отрисовки кадра зависит линейно от количества объектов в сцене. Установлено, что разработанная программная реализация обеспечивает приемлемую производительность при нагрузке в 100 сложных объектов.
\end{enumerate}
